\section{Cancellation at disks and three local bounds}
\label{sec:local}

Throughout this section, $g\in C^\infty(\R^2)$ is real valued,
$h=\nabla g\,dx$, and $u_t=A_tg$. The averages are signed. No
compact-support assumption on $g$ is needed for the local estimates.
Our goal is to control three quantities on a unit square: unfavorable
directional gradients at locally maximizing radii, gradients at band
maxima and endpoints, and differences of averages. The common input is
a decomposition of the gradient into positive directions and error
pieces, with small support and zero total mass below size one. The discount on the latter errors will make the
scale sum in Section~\ref{sec:prices} finite.
We fix
\begin{equation}\label{eq:local-alpha}
\alpha=\frac1{100}.
\end{equation}
An equality of measures on an open set means equality of their
restrictions there. The component measures themselves will always
remain defined on the whole plane.

\begin{lemma}[Local estimates]\label{lem:local}
Let $Q$ be a square of side length one and let
$\Omega=\{y:\dist(y,Q)<10\}$. Suppose that, on $\Omega$,
\begin{equation}\label{eq:local-decomposition}
h=\sum_i v_i\sigma_i+\beta,
\qquad \beta=\sum_d\beta_d.
\end{equation}
The sums are finite. The measures in the sums are finite, defined on
$\R^2$, and absolutely continuous with respect to Lebesgue measure.
Each $\sigma_i$ is nonnegative and each $v_i$ is a constant unit vector.
For each $d$, write $m_d=|\beta_d|(\R^2)$ and assign a number
$s_d\in(0,1]$. If $s_d<1$, assume that
$\beta_d(\R^2)=0$ and that $\beta_d$ is supported in a ball of radius
at most $C_s s_d$, where $C_s$ is fixed. Set
\[
S=\sum_i\sigma_i,
\qquad A=\sum_d s_d^\alpha m_d.
\]
For $x\in Q$, define
\begin{align*}
\mathcal I(x)
&=\{r\in[1,2]:u_r(x)\ge u_t(x)
                 \text{ for every }t>0\text{ with }|t-r|\le1/4\},\\
\mathcal J(x)
&=\{r\in[1,2]:u_r(x)=\max_{1\le t\le2}u_t(x)\}\cup\{1,2\}.
\end{align*}
Then the following estimates hold.
\begin{enumerate}
\item \emph{Unfavorable directional gradients.} If $e$ is a unit vector and $e\cdot v_i\ge\kappa>0$ for all $i$,
then
\begin{equation}\label{eq:local-direction}
\int_Q\sup_{r\in\mathcal I(x)}
                  (-e\cdot\nabla u_r(x))_+\,dx\le C A.
\end{equation}
\item \emph{Band maxima and forced endpoints.} Without a condition on the directions $v_i$,
\begin{equation}\label{eq:local-size}
\int_Q\sup_{r\in\mathcal J(x)}|\nabla u_r(x)|\,dx
    \le C\bigl(A+S(\R^2)\bigr).
\end{equation}
\item \emph{Differences of nearby averages.} We also have
\begin{equation}\label{eq:local-value}
\int_Q\sup_{0<t,t'\le4}|u_t(x)-u_{t'}(x)|\,dx
    \le C\bigl(A+S(\R^2)\bigr).
\end{equation}
\end{enumerate}
Here $C$ depends only on $C_s$, and in the first estimate also on
$\kappa$. A supremum of nonnegative quantities over an empty set is zero.
\end{lemma}

All the quantities in this statement are measurable. To see this,
extend the displayed definitions of $\mathcal I(x)$ and $\mathcal J(x)$
to all $x\in\R^2$. Their graphs are closed in $\R^2\times[1,2]$,
by joint continuity of $u_t(x)$ and continuity of a maximum over a
compact interval. For either of the first two suprema,
its superlevel set at a positive closed threshold is the projection of
a closed set with the radius variable in a compact interval, and is
therefore closed. Restricting these Borel functions to $Q$ proves the
claim also for half-open squares. The last supremum can be taken over positive rational
$t,t'\le4$, by continuity in the radius variables.

The first estimate discounts cancellation without charging the positive
mass $S$: its cone contribution is favorable at every tested radius.
Comparisons of kernels will still require a bound on its weighted mass.
The second estimate pays for
$S$ and also covers the two forced radii, so it can be used when a
radius interval ends. The third controls values rather than gradients;
its factor under rescaling will pay for derivatives of spatial
weights. Section~\ref{sec:prices} gives those scaled applications.

We prove the first two bounds by comparing successively narrower disk
kernels. Before estimating their changes, we explain why a maximizing
radius supplies a useful moment identity and why concentration in a
small cap makes that identity effective.

\subsection{Disk moments and an affine cancellation}

We seek cancellation in a thin shell where two nearby
disks differ. Estimating by the entire gradient mass on that shell
would ignore the maximizing property of the radius. At an interior
maximizing radius, a radial derivative vanishes; a rotational moment
vanishes for every disk. We combine these identities using complex
notation, then preserve them while smoothing the disk boundary.

Variations of optimizing balls also give weighted-gradient cancellation
identities in Luiro's work on the uncentered maximal operator
\cite[Lemmas~2.2 and~2.4]{Luiro2018Radial}. There the balls optimize
averages of $|f|$. We derive the identities needed here directly for
centered signed averages.

We identify the plane with the complex numbers, using the usual positive
orientation. A vector measure is then a complex measure, and
$\overline{y-x}\,dh(y)$ denotes complex multiplication. For every disk,
\begin{equation}\label{eq:local-moment-identities}
\operatorname{Im}\int_{B(x,t)}\overline{y-x}\,dh(y)=0,
\qquad
\partial_tu_t(x)=\frac1{\pi t^3}
 \operatorname{Re}\int_{B(x,t)}\overline{y-x}\,dh(y).
\end{equation}
The first identity follows by integrating the divergence of
$g(y)(y-x)^\perp$: the vector field $(y-x)^\perp$ is divergence free
and tangent to the boundary circle. For the second, differentiate
$u_t(x)=\pi^{-1}\int_{B(0,1)}g(x+tw)\,dw$.
The same change of variables gives
\begin{equation}\label{eq:local-spatial-gradient}
\nabla u_t(x)=\frac{h(B(x,t))}{\pi t^2}.
\end{equation}

\begin{lemma}[Balanced kernels]\label{lem:local-balanced}
Suppose that $r\in[1,2]$ and $0<\delta\le1/4$, and that
\[
u_r(x)\ge u_{r-\delta}(x),\qquad u_r(x)\ge u_{r+\delta}(x).
\]
There is a probability
density $k$ supported on $[r-\delta,r+\delta]$, bounded by $C/\delta$,
such that
\[
K(y)=\int k(t)\mathbf1_{B(x,t)}(y)\,dt
\]
satisfies
\begin{equation}\label{eq:local-zero-moment}
\int K(y)\overline{y-x}\,dh(y)=0.
\end{equation}
Moreover, $0\le K\le1$, $\Lip(K)\le C/\delta$,
$K=1$ on $B(x,r-\delta)$, and $K=0$ outside $B(x,r+\delta)$.
The density may be chosen measurably in $(x,r)$.
\end{lemma}

\begin{proof}
Normalize $t^{-3}$ separately on $[r-\delta,r]$ and
$[r,r+\delta]$ to obtain probability densities $k_-$ and $k_+$.
Let $M_-$ and $M_+$ be the real moments in
\eqref{eq:local-moment-identities} averaged with these densities.
Integration of that identity shows $M_-\ge0$ and $M_+\le0$.
If $M_--M_+>0$, take
\[
k=a k_-+(1-a)k_+,
\qquad a=\frac{-M_+}{M_--M_+};
\]
if both moments are zero, take $k=k_-$. The imaginary moments are
zero separately, so \eqref{eq:local-zero-moment} follows. The
normalizing factors are comparable to $\delta$, because
$3/4\le t\le9/4$. This proves the density bound and all the claims
about $K$. The formula for $a$ is Borel, with the stated convention
at zero, which gives the last assertion.
\end{proof}


The balanced kernels in Lemma~\ref{lem:local-balanced} smooth the disk
boundary without losing its zero moment. Consider two such kernels,
possibly with different widths. Their difference is supported in the
larger shell. The following elementary observation explains what we
want to know about the mass on that shell.

\begin{lemma}[Affine cancellation]\label{lem:local-affine}
Let $x\in\R^2$, $r\in[1,2]$, $0<\delta\le1/4$, and
\[
 D=\{y:\bigl||y-x|-r\bigr|\le\delta\}.
\]
Let $h$ be a complex measure whose variation is finite on $B(x,3)$.
Suppose that $K_0,K_1:\R^2\to[0,1]$ are Borel functions supported in
$B(x,3)$, that $K_1-K_0$ is supported in $D$, and that
\[
 \int K_i(y)\overline{y-x}\,dh(y)=0\qquad(i=0,1).
\]
Set $P_i=\int K_i\,dh$. For every $z\in D$,
\begin{equation}\label{eq:local-affine-identity}
 P_1-P_0=\int
 \left(1-\frac{\overline{y-x}}{\overline{z-x}}\right)
 (K_1-K_0)(y)\,dh(y).
\end{equation}
For every $0<\rho\le1$, this implies
\begin{align}\label{eq:local-affine-bound}
 |P_1-P_0|\le{}&2\rho\int_{D\cap B(z,\rho)}
                    |K_1-K_0|\,d|h|\notag\\
 &+10\int_{D\setminus B(z,\rho)}|K_1-K_0|\,d|h|.
\end{align}
\end{lemma}

\begin{proof}
Since $|z-x|\ge r-\delta\ge3/4$, the denominator is nonzero.
Subtract the difference of the two zero-moment identities, divided by
$\overline{z-x}$, from $P_1-P_0$. This gives
\eqref{eq:local-affine-identity}. Its multiplier is
\[
 W_z(y)=\frac{\overline{z-y}}{\overline{z-x}}.
\]
On $B(z,\rho)$ its absolute value is at most $2\rho$. On $D$,
$|y-z|\le2(r+\delta)\le9/2$ and $|z-x|\ge3/4$, so
$|W_z|\le6\le10$. Splitting the integral into the two regions proves
\eqref{eq:local-affine-bound}.
\end{proof}

Thus concentration near $z$ turns the factor $1$ in a variation
estimate into the small factor $2\rho$. We will call
$D\cap B(z,\rho)$ a cap of the shell. The geometric question is now
concrete: when can all but a controlled amount of shell mass be placed
in one cap?

For the intended use, two positive measures play different roles.
One, denoted by $\nu$, detects whether the shell carries enough error
mass to require this argument. The other, denoted by $S$, describes
positive contributions to the gradient. The conclusion must control
both measures on the shell, but the available hypothesis bounds only
$S$ in the inner disk. A bound involving all of $S(\R^2)$ would lose
that information. The tangent-disk sweep below is what allows the
inner mass bound to suffice.

\subsection{One exceptional set for all radii}

We use Lebesgue measure for the size $|E|$ of a set of centers. The
measures $S$ and $\nu$ in the next lemma are arbitrary finite positive
Borel measures; they may have atoms. No maximizing radius is selected,
and no regularity of these measures is assumed.

\begin{lemma}[Concentration in a thin shell]\label{lem:shell-concentration}
Let $S$ and $\nu$ be finite positive Borel measures on $\R^2$, let
$U>0$ and $T\ge0$, and suppose that
\[
0<\rho\le1,
\qquad 0<\delta\le\frac{\rho^2}{1600}.
\]
Write
\[
D(x,r,\delta)=\{y:\bigl||y-x|-r\bigr|\le\delta\}.
\]
There is a Borel set $E\subset\R^2$ such that
\begin{equation}\label{eq:shell-exception}
|E|\le C\left[
  \left(\frac{\nu(\R^2)}{U}\right)^2\frac{\delta}{\rho}
  +\frac{\nu(\R^2)}{U}\frac{\delta}{\rho^2}
           \left(1+\frac{T}{U}\right)\right],
\end{equation}
and the following implication holds for every $x\notin E$ and every
$r\in[1,2]$: if
\begin{equation}\label{eq:shell-hypotheses}
\nu(D(x,r,\delta))\ge U,
\qquad S(B(x,r-\delta))\le T,
\end{equation}
then there is $z\in D(x,r,\delta)$ for which
\begin{equation}\label{eq:shell-cap}
(S+\nu)\bigl(D(x,r,\delta)\setminus B(z,\rho)\bigr)\le2U.
\end{equation}
The constant in \eqref{eq:shell-exception} is absolute. In particular,
neither the exceptional set nor its bound involves a choice of $r$.
\end{lemma}

\begin{proof}
If the cap conclusion fails, every shell point $z$ sees more than $U$
of either $\nu$ or $S$ outside its cap. The first alternative will
restrict the centers to thin strips. For the second, nested tangent
disks capture that remote $S$-mass during a short parameter interval,
starting with mass at most $T$.
We first prove the needed geometry, then combine the two alternatives
in one Borel majorant that contains no radius parameter.

Put $v=200\delta/\rho^2$, so $v\le1/8$. We first record the geometry
of disks tangent at a fixed point $z$. For a unit vector $\theta$, the
disks
\[
B(z+L\theta,L),\qquad L>0,
\]
increase with $L$. In fact, for $w=y-z$, membership is the inequality
\begin{equation}\label{eq:shell-entry}
|w|^2<2L\,w\cdot\theta.
\end{equation}
When $w\cdot\theta>0$, the entry parameter is
$L_*(y)=|w|^2/(2w\cdot\theta)$.

Fix $x,z$ with $l=|x-z|\in[1/2,3]$ and put
$\theta=(x-z)/l$. If $|y-z|\ge\rho$ and
$\bigl||y-x|-l\bigr|\le2\delta$, then, with
$a=(y-z)\cdot\theta$,
\[
\bigl||y-z|^2-2la\bigr|
 =\bigl||y-x|^2-l^2\bigr|\le13\delta.
\]
The restriction on $\delta$ gives
$a\ge\rho^2/12>0$, and hence
\begin{equation}\label{eq:shell-entry-window}
|L_*(y)-l|\le78\frac{\delta}{\rho^2}<\frac v2.
\end{equation}
Thus all such points enter strictly between $l-v$ and $l+v$.
Conversely, if $|y-z|\ge\rho$ and
$y\in B(z+(l-v)\theta,l-v)$, then
$a>\rho^2/(2(l-v))\ge\rho^2/6$. It follows that
\[
|y-x|^2<l^2-2va
 \le l^2-\frac{200}{3}\delta<(l-2\delta)^2.
\]
We have proved
\begin{equation}\label{eq:shell-earlier-interior}
\{y:|y-z|\ge\rho\}\cap B(z+(l-v)\theta,l-v)
       \subset B(x,l-2\delta).
\end{equation}
Figure~\ref{fig:tangent-disks} illustrates the nested disks and these
two comparisons.

\begin{figure}[tb]
\centering
\begin{tikzpicture}[scale=1.05,>=Stealth]
 \draw[gray!60,thin] (1.55,0) circle[radius=1.55];
 \draw[black,thick] (2,0) circle[radius=2];
 \draw[gray!75,dashed] (2.45,0) circle[radius=2.45];
 \draw[->,gray] (-.35,0)--(5.35,0) node[right] {\(\theta\)};
 \fill (0,0) circle (1.8pt) node[below left] {\(z\)};
 \fill (1.55,0) circle (1.3pt) node[below=5pt] {\(z+(l-v)\theta\)};
 \fill (2,0) circle (1.6pt) node[above=4pt] {\(x=z+l\theta\)};
 \fill (2.45,0) circle (1.3pt);
 \draw[gray] (2.45,-.07)--(5.3,-.85)
   node[right,black] {\(z+(l+v)\theta\)};
 \fill (2,2) circle (2pt) node[above right] {\(y\)};
 \draw[dotted] (0,0)--(2,2);
 \node[align=left,anchor=west] at (5.3,1.65)
   {\(\displaystyle B(z+L\theta,L)\)\\[2pt]increases with \(L\)};
\end{tikzpicture}
\caption{The tangent-disk sweep. All boundary circles pass through \(z\).
A shell point \(y\) with \(|y-z|\ge\rho\) enters between the parameters
\(l-v\) and \(l+v\), with \(v\) proportional to
\(\delta/\rho^2\). Mass outside \(B(z,\rho)\) already present at
\(l-v\) lies in the inner disk \(B(x,l-2\delta)\). When
\(z\in D(x,r,\delta)\), this disk is contained in the controlled disk
\(B(x,r-\delta)\). The drawing is schematic.}
\label{fig:tangent-disks}
\end{figure}


For $z\in\R^2$ and $\theta$ on the unit circle define
\[
H_{z,\theta}(L)=
 \begin{cases}
 S\bigl(\{y:|y-z|\ge\rho\}\cap B(z+L\theta,L)\bigr),&L>0,\\
 0,&L\le0.
 \end{cases}
\]
This is bounded and nondecreasing in $L$. It is jointly Borel in all
its arguments, since it is the integral of the Borel indicator given
by \eqref{eq:shell-entry}, together with the other displayed tests.
For $l=|x-z|\in[1/2,3]$ and $\theta=(x-z)/l$, let $\mathcal C(x,z)$
be the condition
\begin{equation}\label{eq:shell-ray-condition}
H_{z,\theta}(l-v)\le T,
\qquad H_{z,\theta}(l+v)-H_{z,\theta}(l-v)\ge U.
\end{equation}
Define the nonnegative Borel function
\begin{align}
\mathcal M(x)=\frac1U\int_{1/2\le|x-z|\le3}
 \Bigg[&\frac1U\int
 \mathbf1_{\{\,\left||y-x|-|x-z|\right|\le2\delta,
                         \ |y-z|\ge\rho\,\}}\,d\nu(y)
 \notag\\[-2mm]
 &+\mathbf1_{\{\mathcal C(x,z)\}}\Bigg]\,d\nu(z),
 \qquad E=\{x:\mathcal M(x)\ge1\}.
 \label{eq:shell-majorant}
\end{align}

Suppose that \eqref{eq:shell-hypotheses} holds at $(x,r)$ and that
\eqref{eq:shell-cap} fails for every $z\in D=D(x,r,\delta)$.
For any such $z$, either
$\nu(D\setminus B(z,\rho))>U$ or
$S(D\setminus B(z,\rho))>U$.
Also $l=|x-z|\in[1/2,3]$ and every $y\in D$ satisfies
$\bigl||y-x|-l\bigr|\le2\delta$. In the first alternative, the
first term in brackets in \eqref{eq:shell-majorant} exceeds one.
In the second alternative, \eqref{eq:shell-entry-window} gives the
second inequality in \eqref{eq:shell-ray-condition}. The first
inequality follows from \eqref{eq:shell-earlier-interior}, because
$l-2\delta\le r-\delta$. The bracket is therefore at least one for
every $z\in D$. Since $\nu(D)\ge U$, we obtain
$\mathcal M(x)\ge1$. This proves the required implication outside
the Borel set $E$, simultaneously for all $r$.

It remains to estimate $\int\mathcal M$. For fixed $y,z$ with
$|y-z|\ge\rho$, the centers in the first term satisfy $|x-z|\le3$
and
\[
\bigl||x-y|^2-|x-z|^2\bigr|\le13\delta.
\]
The difference of squares is an affine function of $x$ with gradient
of length $2|y-z|$. These centers lie in a strip of width at most
$13\delta/|y-z|$ intersected with $\overline B(z,3)$, and thus have area at
most $C\delta/\rho$. Tonelli's theorem bounds the first contribution
to $\int\mathcal M$ by
$C\nu(\R^2)^2\delta/(U^2\rho)$.

For the second contribution, fix $z,\theta$ and abbreviate
$H=H_{z,\theta}$ and $G=\min(H,T+U)$. On the set of $l$ satisfying
\eqref{eq:shell-ray-condition},
\[
G(l+v)-G(l-v)\ge U.
\]
Because $G$ is bounded and nondecreasing,
\[
\int_{\R}\bigl(G(l+v)-G(l-v)\bigr)\,dl\le2v(T+U).
\]
For completeness, integration over $[-M,M]$ cancels the overlapping
intervals and leaves two intervals of length $2v$; letting $M$ increase
proves the inequality. Thus the permitted $l$ have length at most
$2v(1+T/U)$. Polar coordinates about $z$, with $l\le3$, and another
application of Tonelli's theorem bound the second contribution by
$C\nu(\R^2)v(1+T/U)/U$. Since $|E|\le\int\mathcal M$, these two
estimates prove \eqref{eq:shell-exception}.
\end{proof}


Under its shell and inner-disk mass hypotheses, the shell lemma
supplies a cap for every radius at a center outside one Borel set. The affine identity then suppresses the mass on that
cap. To obtain the discounted estimates, we still have to compare
kernel widths, control the positive interior mass in the directional
argument, and treat radii forced at band endpoints. We do this next;
the value estimate will require a separate kernel.

\subsection{Cancellation of small error pieces}

We now compare kernels at dyadic widths. Put
\begin{equation}\label{eq:local-exponents}
b=2\alpha=\frac1{50},\qquad
\eta=\frac1{10},\qquad \gamma=\frac15,
\qquad \delta_j=2^{-j}.
\end{equation}
Choose a fixed $\delta_*>0$ so small that $\delta_*\le1/4$ and
$\delta\le\delta^{2\gamma}/1600$ whenever $0<\delta\le\delta_*$.

We record estimates valid for any sequence of probability-mixture
kernels $K_j$, centered at the same $x$ and $r\in[1,2]$, with the
support and Lipschitz properties of Lemma~\ref{lem:local-balanced} at
width $\delta_j$. The zero-moment condition is not needed for these
estimates. Write
\begin{equation}\label{eq:local-coarse-measure}
\nu_j=\sum_{s_d\ge\delta_j^2}|\beta_d|,
\qquad p_j=\int K_j\,d\beta,
\qquad D_j=D(x,r,\delta_j).
\end{equation}
Then
\begin{equation}\label{eq:local-basic-cancellation}
\nu_j(\R^2)\le A\delta_j^{-b},
\qquad |p_j|\le C A\delta_j^{-\alpha}.
\end{equation}
Indeed, for $s_d<1$, subtracting the value of $K_j$ at the center of
a supporting ball gives
\[
\left|\int K_j\,d\beta_d\right|
 \le C m_d\min(1,s_d/\delta_j).
\]
For $s_d=1$, the same upper bound holds directly from $|K_j|\le1$.
Now use $\min(1,a)\le a^\alpha$ for $a>0$ and the definition of $A$.
The bound for $\nu_j$ follows directly from $s_d\ge\delta_j^2$.

If $F$ is a scalar or complex test function with
$\Lip(F)\le C/\delta_j$, the pieces with $s_d<\delta_j^2$ satisfy
\begin{equation}\label{eq:local-fine-cancellation}
\left|\sum_{s_d<\delta_j^2}\int F\,d\beta_d\right|
 \le \frac{C}{\delta_j}\sum_{s_d<\delta_j^2}s_dm_d
 \le C A\delta_j^{1-b}.
\end{equation}
Only the Lipschitz bound is used here; constants cancel because each
of these pieces has zero total vector mass. In particular,
\eqref{eq:local-fine-cancellation} applies to $K_{j+1}-K_j$ and to
$\overline{y-x}(K_{j+1}-K_j)$. These tests have the required global
Lipschitz bounds, since their supports are in $B(x,3)$.
Their supports are also in $D_j$. Consequently
\begin{align}
|p_{j+1}-p_j|
&\le C\bigl(\nu_j(D_j)+A\delta_j^{1-b}\bigr),
\notag\\
\left|\int\overline{y-x}(K_{j+1}-K_j)\,d\beta(y)\right|
&\le C\bigl(\nu_j(D_j)+A\delta_j^{1-b}\bigr).
\label{eq:local-increment-tests}
\end{align}

All these kernels are supported in $\Omega$ when $x\in Q$. Thus
\eqref{eq:local-decomposition} can be used in every pairing below,
even if some of its component measures have support outside $\Omega$.
For fixed $x,r$, as $j\to\infty$ the kernels converge off the circle
to $K_\infty=\mathbf1_{B(x,r)}$. Absolute continuity of all the
component measures gives convergence of their integrals by dominated
convergence. The same holds for the moment integrals. In particular, if every
finite-width kernel is balanced, then $K_\infty$ has zero moment as
well. This is a statement at each fixed $(x,r)$: every circle has zero
mass for every absolutely continuous component. There is no exceptional
set depending on the radius to remove. Absolute continuity is used
here, but was not needed in the Borel shell lemma.

\subsection{The directional estimate}

We now prove \eqref{eq:local-direction}. If $A=0$, all the error
measures vanish, and $e\cdot h(B(x,r))\ge0$ for every relevant disk;
the conclusion follows. Assume $A>0$.

We prove a superlevel-set estimate by separating widths at which the
coarse error mass in the shell is small from those at which it is large.
Small-mass widths have directly summable error increments. At the
remaining widths we compare only increases of the negative projection:
the later negative projection controls the interior positive mass,
allowing Lemma~\ref{lem:shell-concentration} to be applied.

Let $\lambda\ge L_0A$, where the fixed constant $L_0$ will be chosen
large enough in the course of the proof. Put $q=A/\lambda$ and choose
$j_0$ so that
\begin{equation}\label{eq:local-start-scale}
q/2<\delta_{j_0}\le q.
\end{equation}
We require $L_0\ge1/\delta_*$, so all widths under consideration are
at most $\delta_*$. At a fixed $x\in Q$ and $r\in\mathcal I(x)$,
use Lemma~\ref{lem:local-balanced} to choose $K_j$ for every
$j\ge j_0$, and write
\[
P_j=\int K_j\,dh,
\qquad N_j=(-e\cdot P_j)_+,
\qquad U_j=\lambda\delta_j^\eta.
\]
Call $j$ active for $(x,r)$ if $\nu_j(D_j)\ge U_j$.
If it is inactive, \eqref{eq:local-increment-tests} and
$1-b>\eta$ give
\begin{equation}\label{eq:local-inactive}
|p_{j+1}-p_j|
+\left|\int\overline{y-x}(K_{j+1}-K_j)\,d\beta(y)\right|
 \le C\lambda\delta_j^\eta.
\end{equation}
Here and below we use $A\le\lambda$.

Choose $C_0$ sufficiently large depending only on $C_s$ and $\kappa$,
and put $T_j=C_0\lambda\delta_j^{-b}$. Apply
Lemma~\ref{lem:shell-concentration} with
\[
\nu=\nu_j,\quad \delta=\delta_j,\quad
\rho=\delta_j^\gamma,\quad U=U_j,\quad T=T_j.
\]
It gives Borel sets $E_j$, independent of $r$, such that
\begin{equation}\label{eq:local-shell-tail}
|E_j|\le C\left[
q^2\delta_j^{14/25}+q\delta_j^{9/25}\right].
\end{equation}
To check the exponents, the first term of
\eqref{eq:shell-exception} gives
$1-\gamma-2b-2\eta=14/25$. The second term is a sum of powers
$1-2\gamma-b-\eta$ and $1-2\gamma-2b-2\eta=9/25$; because
$\delta_j\le1$, both are bounded by a constant times the latter.
Constants in \eqref{eq:local-shell-tail} may depend on the already
fixed $C_0$.

Fix a center outside $\bigcup_{j\ge j_0}E_j$, and keep any
$r\in\mathcal I(x)$ fixed. Positivity in the direction $e$ implies
\begin{equation}\label{eq:local-cone-sign}
e\cdot P_j\ge\kappa\int K_j\,dS+e\cdot p_j,
\qquad N_j\le|p_j|.
\end{equation}
Before the first active index, all the increments of $p_j$ obey
\eqref{eq:local-inactive}. Hence, if $j$ is the first active index,
or if $j=\infty$ and there are no active indices,
\begin{equation}\label{eq:local-first-active}
N_j\le C\left(A\delta_{j_0}^{-\alpha}
                   +\lambda\delta_{j_0}^\eta\right).
\end{equation}
The formula at infinity follows from the dominated-convergence
observation above.

The first active width has therefore been controlled. To continue,
we compare consecutive active widths, allowing arbitrarily many
inactive widths between them. Only an increase of the negative
projection needs a new estimate; that sign will supply the positive
interior-mass bound required by the shell lemma.

Suppose now that $j$ is active and $l>j$ is the next active index;
set $l=\infty$ if there is no next one. The individual bound in
\eqref{eq:local-basic-cancellation} at $j$ and $j+1$, followed by
the inactive increments between $j+1$ and $l$, gives
\begin{equation}\label{eq:local-gap-errors}
|p_j|+|p_l|
 \le C\left(A\delta_j^{-\alpha}
                  +\lambda\delta_j^\eta\right)
 \le C\lambda\delta_j^{-b}.
\end{equation}
We need to estimate a change only when $N_l>N_j$. In that case
$e\cdot P_l<0$, so \eqref{eq:local-cone-sign} gives
\[
\int K_l\,dS\le\kappa^{-1}|p_l|.
\]
Since $K_l=1$ on $B(x,r-\delta_j)$, this proves
$S(B(x,r-\delta_j))\le T_j$ after fixing $C_0$ large enough to
dominate the constant just obtained. The shell lemma therefore
supplies a point $z\in D_j$ such that
\begin{equation}\label{eq:local-chosen-cap}
(S+\nu_j)\bigl(D_j\setminus B(z,\delta_j^\gamma)\bigr)
 \le2U_j.
\end{equation}
This point is used only at the current center and radius; no
measurable selection of $z$ is required.

To estimate the positive mass on the cap, we need a weighted bound
for both kernels being compared. The sign at $l$ supplies one, but
$e\cdot P_j>0$ would give no corresponding bound for $K_j$.
We therefore replace $K_j$, when necessary, by a kernel with zero
projection and the same negative part.
If $e\cdot P_j\le0$, set $K_*=K_j$. Otherwise, interpolate between
$K_j$ and $K_l$ to obtain a convex combination $K_*$ with
$e\cdot P_*=0$, where $P_* =\int K_*\,dh$. This is possible because
$e\cdot P_l<0$. In either case,
\[
N_*:=(-e\cdot P_*)_+=N_j,
\qquad e\cdot P_*\le0.
\]
Writing $p_* =\int K_*\,d\beta$ and using
\eqref{eq:local-gap-errors} and \eqref{eq:local-cone-sign}, we get
\begin{equation}\label{eq:local-weighted-good}
\int(K_l+K_*)\,dS\le C\lambda\delta_j^{-b}.
\end{equation}

All of $K_j,K_l,K_*$ have zero moment, including when $l=\infty$
by the preceding sharp-kernel limit. Since $|z-x|\ge1/2$, define
\[
W(y)=1-\frac{\overline{y-x}}{\overline{z-x}}.
\]
The affine identity in Lemma~\ref{lem:local-affine} gives
\begin{equation}\label{eq:local-corrected-difference}
P_l-P_* =\int W(y)(K_l-K_*)(y)\,dh(y).
\end{equation}
The difference of kernels is supported on $D_j$, and
\[
|W|\le C\delta_j^\gamma\quad\hbox{on }B(z,\delta_j^\gamma),
\qquad |W|\le C\quad\hbox{on }D_j.
\]
The contribution of $\sum_i v_i\sigma_i$ to
\eqref{eq:local-corrected-difference} is therefore at most
\begin{equation}\label{eq:local-good-cap-cost}
C\lambda\delta_j^{\gamma-b}+C U_j,
\end{equation}
by \eqref{eq:local-weighted-good},
$|K_l-K_*|\le K_l+K_*$, and \eqref{eq:local-chosen-cap}.

For the error measures, $K_l-K_*$ is a scalar multiple in $[0,1]$
of $K_l-K_j$. Keep $W$ fixed and telescope
\begin{equation}\label{eq:local-error-telescoping}
\int W(K_l-K_j)\,d\beta
 =\sum_{i=j}^{l-1}\int W(K_{i+1}-K_i)\,d\beta.
\end{equation}
For $l=\infty$ this equality means the limit of the finite sums.
We never apply a Lipschitz cancellation bound directly to the sharp
kernel $K_\infty$ or to its interpolant: only the finite consecutive
differences in this sum need such a bound.
On the first step, the pieces contributing to $\nu_j$ cost at most
$C\delta_j^\gamma\nu_j(\R^2)+CU_j$ by
\eqref{eq:local-chosen-cap}. The remaining pieces cost
$CA\delta_j^{1-b}$ by \eqref{eq:local-fine-cancellation}, since
$W(K_{j+1}-K_j)$ has Lipschitz constant at most $C/\delta_j$.
For $j<i<l$, use
\[
\int W(K_{i+1}-K_i)\,d\beta
 =p_{i+1}-p_i-\frac1{\overline{z-x}}
       \int\overline{y-x}(K_{i+1}-K_i)\,d\beta(y).
\]
These are inactive steps, so their absolute values sum to at most
$C\lambda\delta_j^\eta$ by \eqref{eq:local-inactive}.
Combining the estimates gives
\[
\left|\int W(K_l-K_j)\,d\beta\right|
 \le C\left(A\delta_j^{\gamma-b}
       +A\delta_j^{1-b}+\lambda\delta_j^\eta\right).
\]
Together with \eqref{eq:local-good-cap-cost}, and using
$\gamma-b=9/50>\eta$, this yields
\begin{equation}\label{eq:local-active-increment}
N_l\le N_j+|P_l-P_*|\le N_j+C\lambda\delta_j^\eta.
\end{equation}
The same inequality is automatic when $N_l\le N_j$.

Sum \eqref{eq:local-active-increment} along the active indices. If
there are finitely many, include the last comparison with $l=\infty$;
if there are infinitely many, pass to the limit along them. Dominated
convergence and \eqref{eq:local-first-active} give
\begin{equation}\label{eq:local-final-projection}
(-e\cdot h(B(x,r)))_+
 \le C\left(A\delta_{j_0}^{-\alpha}
                   +\lambda\delta_{j_0}^\eta\right)
 \le C\lambda\left(q^{1-\alpha}+q^\eta\right).
\end{equation}
This holds for every $r\in\mathcal I(x)$ at the fixed center outside
the exceptional sets. Increase $L_0$, depending only on $C_s$ and
$\kappa$, so that the last bound is less than $\lambda$ whenever
$q\le L_0^{-1}$.

Let $Y(x)=\sup_{r\in\mathcal I(x)}(-e\cdot h(B(x,r)))_+$.
Summing \eqref{eq:local-shell-tail} and using
\eqref{eq:local-start-scale}, we have proved
\begin{equation}\label{eq:local-direction-tail}
|\{x\in Q:Y(x)>\lambda\}|
 \le C\left[
       (A/\lambda)^{64/25}+(A/\lambda)^{34/25}\right]
 \qquad(\lambda\ge L_0A).
\end{equation}
Both powers exceed one. The layer-cake formula, the trivial bound
$|Q|=1$ for smaller thresholds, and the substitution
$\lambda=As$ give
\[
\int_QY\,dx
 \le L_0A+C A\int_{L_0}^{\infty}
                  (s^{-64/25}+s^{-34/25})\,ds\le C A.
\]
Finally, \eqref{eq:local-spatial-gradient} and $r\ge1$ prove
\eqref{eq:local-direction}.

\subsection{The size estimate, including band endpoints}

To prove \eqref{eq:local-size}, place the measures $v_i\sigma_i$
among the error measures with assigned size one. The augmented price is
$A_{\mathrm{tot}}=A+S(\R^2)$, where $A$ and $S$ refer to the original
decomposition. The augmented positive part is zero; in this proof,
$\beta$ and $\nu_j$ refer to the augmented error family.
If $A_{\mathrm{tot}}=0$, the result is immediate.

Take $\lambda\ge L_0A_{\mathrm{tot}}$, put $q=A_{\mathrm{tot}}/\lambda$, and choose $j_0$ as in
\eqref{eq:local-start-scale}. Here $L_0$ will depend only on $C_s$. For $x\in Q$ and
$r\in\mathcal J(x)$, choose the balanced kernel $K_j$ when
$\dist(r,\{1,2\})\ge\delta_j$. Such an $r$ is an actual
maximizer on $[1,2]$, so the two neighboring radii needed in
Lemma~\ref{lem:local-balanced} are admissible. Otherwise choose the
uniform probability mixture over $[r-\delta_j,r+\delta_j]$.
The latter kernel need not have zero moment, but has all the same
support, positivity and Lipschitz bounds. Thus
\eqref{eq:local-basic-cancellation}--\eqref{eq:local-increment-tests}
remain valid with $A_{\mathrm{tot}}$ in place of $A$. Write $P_j=\int K_j\,dh=p_j$, and again call $j$
active when $\nu_j(D_j)\ge U_j=\lambda\delta_j^\eta$.

There is an additional Borel exceptional set for each width:
\begin{equation}\label{eq:local-endpoint-exception}
F_j=\bigcup_{a\in\{1,2\}}
 \left\{x:\nu_j\{y:\bigl||y-x|-a\bigr|\le2\delta_j\}
                         \ge U_j\right\}.
\end{equation}
For each fixed $y$, the areas of these two annuli of centers sum to
$24\pi\delta_j$. Tonelli's theorem and Markov's inequality give
\begin{equation}\label{eq:local-endpoint-area}
|F_j|\le24\pi\delta_j\frac{\nu_j(\R^2)}{U_j}
 \le C(A_{\mathrm{tot}}/\lambda)\delta_j^{22/25}.
\end{equation}
If $\dist(r,\{1,2\})<\delta_j$, the shell $D_j$ lies in one of
the annuli in \eqref{eq:local-endpoint-exception}. Thus, outside
$F_j$, an active index must satisfy
$\dist(r,\{1,2\})\ge\delta_j$, and both $K_j$ and $K_{j+1}$
have zero moment. In particular, the forced radii $r=1,2$ are
covered even if they are not maximizers.

Apply Lemma~\ref{lem:shell-concentration} with $S=0$, $T=0$,
$\nu=\nu_j$, $\delta=\delta_j$, $\rho=\delta_j^\gamma$, and $U=U_j$.
It supplies Borel exceptional sets $E_j$, independent of $r$, with bound
$C[q^2\delta_j^{14/25}+q\delta_j^{12/25}]$, which is at most the
right side of \eqref{eq:local-shell-tail}. Its constants have no cone
parameter: the positive measure and its interior-mass threshold are
both zero.
At an active step
outside $E_j\cup F_j$, its interior-mass hypothesis is automatic,
and it gives a point $z\in D_j$ with
$\nu_j(D_j\setminus B(z,\delta_j^\gamma))\le2U_j$.
The zero moments allow the correction
\[
P_{j+1}-P_j=\int
 \left(1-\frac{\overline{y-x}}{\overline{z-x}}\right)
 (K_{j+1}-K_j)(y)\,d\beta(y).
\]
The coarse pieces cost at most
$C A_{\mathrm{tot}}\delta_j^{\gamma-b}+CU_j$, and the fine pieces cost at most
$C A_{\mathrm{tot}}\delta_j^{1-b}$, exactly as above. Hence
\[
|P_{j+1}-P_j|\le C\lambda\delta_j^\eta.
\]
At an inactive step the same inequality follows directly from
\eqref{eq:local-increment-tests}, without using a moment identity.
Starting with $|P_{j_0}|\le CA_{\mathrm{tot}}\delta_{j_0}^{-\alpha}$ and summing
all increments, we obtain
\[
|h(B(x,r))|\le C\lambda(q^{1-\alpha}+q^\eta)<\lambda
\]
after increasing $L_0$. This holds simultaneously for all
$r\in\mathcal J(x)$ outside $\bigcup_{j\ge j_0}(E_j\cup F_j)$.
The measure bound is therefore
\begin{align}
\left|\left\{x\in Q:
       \sup_{r\in\mathcal J(x)}|h(B(x,r))|>\lambda\right\}\right|
 \le C\bigl[&(A_{\mathrm{tot}}/\lambda)^{64/25}
              +(A_{\mathrm{tot}}/\lambda)^{34/25}
 \notag\\[-1mm]
              &+(A_{\mathrm{tot}}/\lambda)^{47/25}\bigr].
\label{eq:local-size-tail}
\end{align}
Every exponent is larger than one. The same layer-cake computation
as before gives the unnormalized bound by $CA_{\mathrm{tot}}$. Since
$A_{\mathrm{tot}}=A+S(\R^2)$,
\eqref{eq:local-spatial-gradient} proves \eqref{eq:local-size}.

\subsection{The value-oscillation estimate}

This estimate involves all small radii and does not use a maximizing
property. We express $g$ minus a fixed mean through its gradient,
prove an $L^{3/2}$ cancellation bound for the resulting kernel, and
then apply the strong maximal inequality.

We finish the proof of Lemma~\ref{lem:local} by establishing
\eqref{eq:local-value}. Choose a nonnegative smooth probability
density $\psi$ supported in the interior of $Q$, with uniformly
bounded derivatives after translating and rotating a fixed choice,
and put
\[
c=\int g(z)\psi(z)\,dz,
\qquad E_Q=\{x:\dist(x,Q)<4\}.
\]
For $x\in E_Q$, integrate $\nabla g$ along the segments from $x$
to points of $\supp\psi$ and change variables to obtain
\begin{equation}\label{eq:local-value-kernel}
g(x)-c=\int L(x,y)\cdot dh(y),
\qquad
L(x,y)=-(y-x)\int_0^1
              \psi\bigl(x+(y-x)/t\bigr)t^{-3}\,dt.
\end{equation}
Indeed, this is the fundamental theorem of calculus in
$g(x)-g(z)=-\int_0^1(z-x)\cdot\nabla g(x+t(z-x))\,dt$,
followed by $y=x+t(z-x)$.

The displayed formula defines $L(x,y)$ for every $y\ne x$ in
the plane, not only for $y\in\Omega$. This global definition will
allow cancellation against a support center outside $\Omega$.
For each fixed $x\in E_Q$, the kernel is supported on segments
joining $x$ to $\supp\psi$, and hence is supported in $\Omega$.
It is smooth in $y\ne x$, with the uniform bounds
\begin{equation}\label{eq:local-value-kernel-bounds}
|L(x,y)|\le\frac{C}{|x-y|},
\qquad |\nabla_yL(x,y)|\le\frac{C}{|x-y|^2}.
\end{equation}
To verify them, if the integrand or one of its derivatives is
nonzero, then $|y-x|\le Ct$, since $x\in E_Q$ and
$x+(y-x)/t$ lies in a fixed bounded neighborhood of $Q$.
The first bound follows by integrating $t^{-3}$ from
$|y-x|/C$ to one and multiplying by $|y-x|$. Differentiation in
$y$ produces either the same integral or
$|y-x|\int t^{-4}\,dt$, which gives the second bound.

Fix $p=3/2$, so $2/p-1=1/3>\alpha$. The first bound in
\eqref{eq:local-value-kernel-bounds} implies
\begin{equation}\label{eq:local-kernel-lp}
\sup_{y\in\R^2}\|L(\cdot,y)\|_{L^p(E_Q)}\le C.
\end{equation}
This follows by integrating $|x-y|^{-p}$ on a bounded set;
$p<2$ makes the singularity integrable, uniformly in $y$.

Consider a piece $\beta_d$ with $s_d<1$, supported in
$B(z,a)$, where we may take $a=C_s s_d$ after increasing $C_s$
to at least one. On $E_Q\cap B(z,2a)$, the first kernel bound and
Minkowski's inequality give
\[
\left\|\int L(\cdot,y)\cdot d\beta_d(y)\right\|_
                       {L^p(E_Q\cap B(z,2a))}
 \le C a^{2/p-1}m_d.
\]
In fact, for $y\in B(z,a)$ the relevant $x$ lie within $3a$ of
$y$, and the $L^p$ norm of $|x-y|^{-1}$ there is
$C a^{2/p-1}$.
For $|x-z|>2a$, use the zero total vector mass to subtract
$L(x,z)$. The derivative bound in
\eqref{eq:local-value-kernel-bounds}, applied on the segment from
$z$ to $y$, gives
\[
\left|\int L(x,y)\cdot d\beta_d(y)\right|
 \le C a m_d|x-z|^{-2}.
\]
Its $L^p$ norm on $\{|x-z|>2a\}$ is at most
$C a^{2/p-1}m_d$, since $p>1$. Combining the two regions yields
\begin{equation}\label{eq:local-value-piece}
\left\|\int L(\cdot,y)\cdot d\beta_d(y)\right\|_{L^p(E_Q)}
 \le C s_d^{1/3}m_d\le C s_d^\alpha m_d.
\end{equation}
The constant may depend on $C_s$. For a piece of assigned size
one, and for the positive measures, use
\eqref{eq:local-kernel-lp} directly. Minkowski's inequality and the
finite decomposition therefore give
\begin{equation}\label{eq:local-value-lp}
\|g-c\|_{L^{3/2}(E_Q)}\le C\bigl(A+S(\R^2)\bigr).
\end{equation}
These same estimates show that the integrals obtained by inserting
\eqref{eq:local-decomposition} into \eqref{eq:local-value-kernel}
exist for almost every $x$; thus the use of the potentially singular
kernel is justified.

Let $f_Q=(g-c)\mathbf1_{E_Q}$, extended by zero to the plane.
Every disk centered at $x\in Q$ of radius at most four is contained
in $E_Q$. Consequently
\[
\sup_{0<t,t'\le4}|u_t(x)-u_{t'}(x)|\le2Mf_Q(x),
\]
where the maximal function on the right uses absolute values.
The classical strong $L^{3/2}$ maximal inequality
\cite[(1.2)]{Kinnunen1997}, followed by H\"older's inequality on the unit square and
\eqref{eq:local-value-lp}, proves \eqref{eq:local-value}. This
completes the proof of Lemma~\ref{lem:local}.
